\ChapterImageStar[cap:trabajos-futuros]{Trabajos Futuros}{./images/fondo.png}\label{cap:trabajos-futuros}
\mbox{}\\

El desarrollo de la notación \NDITI~y su evaluación inicial han permitido establecer una base conceptual y metodológica sólida para el modelado de infraestructura de TI en el contexto académico. No obstante, los resultados obtenidos abren diversas líneas de trabajo futuro orientadas a fortalecer su aplicabilidad, validación y adopción institucional.

\section{Consolidación de la notación en herramientas de software}
\noindent

Uno de los pasos más relevantes para la evolución de la propuesta consiste en la implementación de la notación \NDITI~en una herramienta de software que permita su uso práctico en entornos académicos y, potencialmente, profesionales.

\textbf{Integración en herramientas de modelado}: Se plantea la extensión de herramientas existentes (como editores compatibles con ArchiMate o plataformas de modelado visual) mediante la incorporación de los elementos y reglas de la notación propuesta. Esto permitiría validar su uso en escenarios reales de modelado.

\textbf{Desarrollo de un editor propio}: Como alternativa, se propone el diseño de una herramienta ligera enfocada específicamente en \NDITI, que facilite la creación de diagramas, reduzca la complejidad de uso y garantice la correcta aplicación de la notación.

\textbf{Estandarización interna}: La disponibilidad de una herramienta facilitaría la adopción progresiva de la notación dentro de la Universidad del Quindío, promoviendo su uso como estándar en asignaturas relacionadas con infraestructura de TI.

\section{Evaluaciones académicas en contexto real}
\noindent

La validación realizada en este trabajo constituye un primer acercamiento al uso de la notación; sin embargo, es necesario ampliar su evaluación en escenarios académicos reales.

\textbf{Aplicación en asignaturas}: Se propone implementar la notación \NDITI~en cursos del programa de Ingeniería de Sistemas y Computación, permitiendo que los docentes la utilicen como herramienta de enseñanza en actividades de modelado.

\textbf{Evaluación con docentes}: La participación de profesores en el uso de la notación permitirá identificar fortalezas y limitaciones desde una perspectiva pedagógica, complementando la evaluación realizada con estudiantes.

\textbf{Estudios longitudinales}: La aplicación de la notación a lo largo de varios periodos académicos permitiría analizar su impacto en el aprendizaje, la claridad conceptual y la capacidad de abstracción de los estudiantes.

\section{Evolución y refinamiento de la notación}
\noindent

\textbf{Ajuste de elementos y relaciones}: A partir de futuras evaluaciones, se podrán identificar mejoras en la definición semántica de los elementos, así como en las reglas de relación entre capas.

\textbf{Ampliación de cobertura}: Se propone extender la notación para incluir nuevos dominios relevantes, como aspectos avanzados de seguridad, gobierno de TI o integración con arquitecturas emergentes.

\textbf{Optimización de la complejidad}: Uno de los retos identificados es equilibrar la expresividad con la facilidad de uso, por lo que futuras iteraciones deberán buscar simplificar la notación sin perder capacidad de representación.

\section{Comparación con otras notaciones}
\noindent

\textbf{Estudios comparativos}: Se plantea realizar evaluaciones controladas que comparen el uso de \NDITI~con notaciones como UML, BPMN y ArchiMate en tareas específicas de modelado de infraestructura.

\textbf{Medición de desempeño}: Estos estudios permitirían analizar métricas como tiempo de modelado, tasa de errores, nivel de comprensión y calidad de los diagramas generados.

\section{Adopción institucional y transferencia de conocimiento}
\noindent

\textbf{Materiales educativos}: Desarrollo de guías, ejemplos, plantillas y casos de estudio que faciliten la enseñanza y aprendizaje de la notación.

\textbf{Capacitación docente}: Implementación de talleres o espacios de formación para profesores del programa, orientados al uso efectivo de la notación en sus cursos.

\textbf{Estandarización curricular}: A mediano plazo, se plantea la posibilidad de incorporar \NDITI~como parte de los lineamientos de modelado en el área de infraestructura de TI dentro del currículo académico.

\section{Impacto a largo plazo}
\noindent

Las líneas de trabajo futuro identificadas tienen el potencial de consolidar a \NDITI~como una herramienta relevante para el modelado de infraestructura de TI en el ámbito académico. Su evolución hacia una solución soportada por software, junto con su validación en escenarios reales de enseñanza, puede contribuir a mejorar la comprensión de conceptos complejos y a establecer un lenguaje común entre docentes y estudiantes.

El avance en estas direcciones permitirá fortalecer la coherencia en la representación de sistemas de infraestructura, apoyar los procesos formativos del programa de Ingeniería de Sistemas y Computación y aportar al desarrollo de prácticas más estructuradas en el modelado de TI dentro de la Universidad del Quindío.